\section*{HISTORICAL NOTE}
\phantomsection
\addcontentsline{toc}{section}{Historical Note}

(Numbers in brackets refer to the bibliography at the end of this note.)

The main results of the theory of subgroups and quotient groups of the additive groups \(R^{n}\) have been known since the end of the last century. Many questions of arithmetic and analysis led mathematicians to investigate the structure of subgroups of \(R^{n}\) generated by a finite number of points. Thus Lagrange, while developing the theory of ``continued fractions'', showed in passing that, for every real number \(\theta\) , there exist integers m, n, not both zero, such that \(m - n\theta\) is arbitrarily small ({[}1{]}, vol.~7, p.~27). In 1835 Jacobi, motivated by his researches on periodic analytic functions of several complex variables, showed that if x, y, z are three vectors of \(R^{2}\) , there exist integers m, n, p, not all zero, which make the vector \(mx + ny + pz\) arbitrarily small ({[}2{]}, vol.~2, p.~25). A little later Dirichlet, in the course of his work on the theory of algebraic numbers, discovered his famous ``pigeon-hole principle'' (Schubfachprinzip) (cf.~§ I, Exercises 10 and 11), by means of which he showed that p forms \(\alpha_{i1}m_{1} + \alpha_{i2}m_{2} + \cdots + \alpha_{in}m_{n} - q_{i}\) ( \(1 \leqslant i \leqslant p\) ), where the \(\alpha_{ij}\) are arbitrary real numbers and the \(m_{j}\) and the \(q_{i}\) are integers (not all zero), can simultaneously be made arbitrarily small ({[}3{]}, vol.~I, p.~635). By an entirely different method Hermite arrived at the same result in 1850, for forms of the particular type \(m\theta_{i} - q_{i}\) ( \(1 \leqslant i \leqslant p\) ) ({[}4{]}, vol.~I, p.~105). Finally, in 1884, Kronecker proved the general result stated in Proposition 7 of § I ({[}5{]}, vol.~31, p.~47).

Of course these results were independent of the general theory of locally compact abelian groups, which is of recent origin (see the Historical Note to Chapter III); but this latter theory, particularly the theory of duality (*), has shed a new light on these old results, mainly by bringing out the fundamental concept of associated subgroups. The exposition in the text is based on these ideas (*).

The point of view we have taken in this chapter is purely qualitative: that is to say, we have established the existence of linear combinations of p points, with integer coefficients, which approximate arbitrarily closely to a given point (which may possibly have to satisfy certain conditions); but one may also ask whether there exist relations between the accuracy of the approximation and the magnitude of the coefficients in the linear combinations which give the approximation; this is the point of view of the quantitative theory of ``Diophantine approximation'', and the point of view of all the authors quoted in the first paragraph. Over the last hundred years these questions have been the object of many and diverse investigations, rich in applications to the theory of numbers; to trace their development in this context would take us far outside our present scope, and we shall therefore do no more than refer the reader who wishes to go into these theories to the fundamental writings of Minkowski {[}6{]} and H. Weyl {[}7{]}, which are the origin of an abundant literature ( \(^{**}\) ).

